% SEGMENT OLP-0127-S01
% Part: first-order-logic
% Chapter: completeness
% Section: introduction

\documentclass[../../../include/open-logic-section]{subfiles}

\begin{document}

\iftag{FOL}
      {\olfileid{fol}{com}{int}}
      {\olfileid{pl}{com}{int}}

\olsection{அறிமுகம்}

% SEGMENT OLP-0127-S02
நிறைவு தேற்றம் தருக்கவியலின் மிக அடிப்படையான முடிவுகளில் ஒன்று. அதற்கு
இரண்டு வடிவங்கள் உள்ளன; அவை சமானமானவை என்பதை நாம் நிறுவுவோம். அதன் முதல்
வடிவம் பொருண்மைப் பின்விளைவுக்கும் நமது !!{derivation} அமைப்புக்கும் இடையிலான
அடிப்படை உறவைக் கூறுகிறது: !!a{sentence}~$!A$ ஆனது $\Gamma$ என்ற சில
!!{sentence}s இலிருந்து பின்வருமெனில், $\Gamma \Proves !A$ என்பதை நிறுவும்
!!a{derivation} உண்டு. எனவே, உண்மையில் பின்வராதவற்றை நிறுவாமல் இருக்கும்போது
எவ்வளவு வலிமையாக இருக்க முடியுமோ அவ்வளவு வலிமையானதாக அந்த !!{derivation}
அமைப்பு உள்ளது.

% SEGMENT OLP-0127-S03
அதன் இரண்டாம் வடிவத்தை ஒரு மாதிரி இருப்பு முடிவாகக் கூறலாம்: முரணற்ற
!!{sentence}s கொண்ட ஒவ்வொரு கணமும் நிறைவுறத்தக்கது. முரணின்மை என்பது ஒரு
நிறுவல்-கோட்பாட்டுக் கருத்து: சில வகையான !!{derivation}s ஐ உருவாக்க நமது
!!{derivation} அமைப்பால் இயலாது என்று அது கூறுகிறது. ஆனால்~$\Gamma$ இலிருந்து
ஒரு குறிப்பிட்ட வகை !!{derivation}s இல்லை என்பதால் மட்டுமே,
\iftag{FOL}{!!a{structure}~$\Struct{M}$}{!!{valuation}{~$\pAssign{v}$}
ஒன்று இருந்து $\iftag{FOL}{\Sat{M}}{\pSat{v}}{\Gamma}$} என்பதற்கு என்ன
உத்தரவாதம்? நிறைவு தேற்றம் முதன்முதலில் நிறுவப்படுவதற்கு முன்பே---உண்மையில்,
நாம் இப்போது பயன்படுத்தும் !!{derivation} அமைப்புகள் உருவாவதற்கு முன்பே---ஒரு
கணிதக் கோட்பாட்டின் முரணின்மை, அது பேசும் பொருள்கள் இருப்பதை உறுதிப்படுத்தும்
என்று சிறந்த செருமானியக் கணிதவியலாளர் டேவிட் ஹில்பர்ட் கருதினார். கோட்லோப்
ஃபிரேகேவுக்கு எழுதிய கடிதத்தில் அவர் இதைப் பின்வருமாறு கூறினார்:
\begin{quote}
  தன்னிச்சையாகத் தரப்பட்ட அடிகோள்கள், அவற்றின் எல்லாப் பின்விளைவுகளோடும்,
  ஒன்றுக்கொன்று முரண்படவில்லையெனில் அவை மெய்யானவை; அந்த அடிகோள்களால்
  வரையறுக்கப்படும் பொருள்கள் இருக்கின்றன. எனக்கு இதுவே மெய்மைக்கும்
  இருப்புக்கும் அளவுகோல்.
\end{quote}
ஃபிரேகே இதைக் கடுமையாக மறுத்தார். அடிகோள்கள் முரணற்றவையாக இருந்தால்,
அவை அனைத்தையும் மெய்யாக்கும் \emph{ஏதோவொரு}
\iftag{FOL}{!!{structure}}{!!{valuation}} இருக்கிறது என்ற பொருளிலாவது
ஹில்பர்ட் சரியே என்பதை நிறைவு தேற்றத்தின் இரண்டாம் வடிவம் காட்டுகிறது.

% SEGMENT OLP-0127-S04
நிறைவு தேற்றம்---இன்னும் துல்லியமாகச் சொன்னால் அதன் நிறுவல்---முக்கியமானதற்கு
இவை மட்டுமே காரணங்கள் அல்ல. அதற்கு முக்கியமான பல துணைவிளைவுகள் உள்ளன;
அவற்றில் சிலவற்றைத் தனித்தனியாக விவாதிப்போம். எடுத்துக்காட்டாக,
$\Gamma \Proves !A$ என்பதைக் காட்டும் ஒவ்வொரு !!{derivation} உம்
இறுதிக்குட்பட்டது; எனவே அது~$\Gamma$ விலுள்ள இறுதிக்குட்பட்ட எண்ணிக்கையிலான
!!{sentence}s ஐ மட்டுமே பயன்படுத்த முடியும். ஆதலால் $!A$ ஆனது~$\Gamma$ வின்
பின்விளைவு எனில், அது ஏற்கெனவே~$\Gamma$ வின் ஓர் இறுதிக்குட்பட்ட உட்கணத்தின்
பின்விளைவு என்பது நிறைவு தேற்றத்திலிருந்து கிடைக்கிறது. இதை
\emph{இறுக்கத்தன்மை} என்கிறோம். சமானமாக, $\Gamma$ வின் ஒவ்வோர்
இறுதிக்குட்பட்ட உட்கணமும் முரணற்றது எனில், $\Gamma$ உம் முரணற்றதாக இருக்க
வேண்டும்.

% SEGMENT OLP-0127-S05
!!{derivation}s வழியாகச் சுற்றிச் சென்று நிறைவு தேற்றத்திலிருந்து
இறுக்கத்தன்மைத் தேற்றத்தைப் பெறலாம். அதேவேளையில், நிறைவு தேற்றத்தின்
\emph{நிறுவலைப்} பயன்படுத்தி அதை நேரடியாகவும் நிறுவலாம். ஏனெனில் ஒரு
குறிப்பிட்ட பண்பைக்---முரணின்மையைக்---கொண்ட !!{sentence}s கணத்தை அந்த
நிறுவல் எடுத்துக்கொண்டு, அதிலிருந்து சில பண்புகளைக் கொண்ட !!a{structure}
ஐக் கட்டமைக்கிறது (இங்கு அந்தக் கட்டமைப்பு கணத்தை நிறைவுறுத்துகிறது).
முரணற்ற !!{sentence}s கணங்களுக்குப் பதிலாக ``ஒவ்வோர் இறுதிக்குட்பட்ட
உட்கணமும் நிறைவுறத்தக்க'' கணங்களிலிருந்து தொடங்கினால், ஏறத்தாழ அதே
கட்டமைப்பைப் பயன்படுத்தி இறுக்கத்தன்மையை நேரடியாக நிறுவலாம்.
\iftag{FOL}{இந்தக் கட்டமைப்பு வேறு துணைவிளைவுகளையும் தருகிறது. எடுத்துக்காட்டாக,
  நிறைவுறத்தக்க !!{sentence}s கொண்ட ஒவ்வொரு கணத்துக்கும் இறுதிக்குட்பட்ட அல்லது
  !!{denumerable}s மாதிரி ஒன்று உள்ளது. (இந்த முடிவை
  லோவென்ஹைம்--ஸ்கோலம் தேற்றம் என்கிறோம்.) பொதுவாக, !!{sentence}s கணங்களிலிருந்து
  !!{structure}s ஐக் கட்டமைக்கும் முறை தருக்கவியலில் அடிக்கடியும், சிலவேளைகளில்
  மெய்யியலிலும்கூடப் பயன்படுத்தப்படுகிறது.}{}

\end{document}
